\documentclass[../../main.tex]{subfiles}
\begin{document}
\section{Deterministic moment-map/Haar route: audited frontier}
\label{sec:moment-map-cmh}

This section records a second, deterministic use of the moment map. Its aim is to extend
Theorem~\ref{thm:letwin-moment-map} from constant matrices to multiplier fields selected by an
arbitrary test function. The route is live but unproved. Displayed formal identities below are
not assigned R2 status; complete conventions, proofs, retractions, and regression models are kept
in \texttt{research/kls/routes/moment-map-cmh/}.

To avoid a collision with the stochastic quantities $H_t,S_t,K_t$ and $A_t$ used earlier, all
objects in this section are stationary moment-map objects: $H=D^2\phi$ is the Hessian metric,
$N$ is a weighted elliptic operator, and $K_M$ is a compressed multiplier.

\begin{program}[Deterministic CMH route]
\label{prog:cmh-route}
First make the CMH endpoint and its KLS reduction precise. Then derive the invariant multiplier
lift, control the resulting square-root commutators with the full positive reservoir, and sum the
complete Haar tree without nodewise positivity or duplicated slack.
\end{program}

\subsection{Endpoint and first audit gate}

The originating exploration proposes a covariance--moment--Hessian estimate of the schematic
form
\begin{equation}
  \norm{\Sigma^{-1/2}H\nabla g}_2^2
  \le4\norm{-Lg}_2^2.
  \label{eq:cmh4-schema}
\end{equation}
It reports that a divergence-duality argument would then give $\CP(\mu)\le4$, the sharp plausible
constant because a standard centered one-sided exponential has Poincar\'e constant $4$.

\begin{question}[CMH normalization and endpoint reduction]
\label{q:cmh-normalization}
Define $\Sigma$, $L$, the underlying $L^2$ space, the admissible class of $g$, and every inverse
and closure in \eqref{eq:cmh4-schema}. Then prove, uniformly through regular moment-map
approximation, that the resulting statement implies $\CP(\mu)\le4$ for every isotropic
log-concave $\mu$.
\end{question}

This question is logically prior to the commutator calculation. Until it is answered,
CMH(4) is a route schema rather than a certified conditional theorem.

\subsection{Haar compression and the commutator error}

On the mean-zero subspace, formally set
\[
  N=D_\nu^*H^{-1}D,
  \qquad
  Q_M=D_\nu^*H^{-1/2}MH^{-1/2}D,
\]
and
\[
  R=H^{-1/2}DN^{-1/2},
  \qquad K_M=R^*MR.
\]
On a common core, $R^*R=I$ and
\[
  Q_M=N^{1/2}K_MN^{1/2}.
\]
For normalized Haar contrast multipliers $M_S$, the formal error is therefore
\begin{equation}
  e_S(u)=Q_{M_S}u-K_{M_S}Nu
  =[N^{1/2},K_{M_S}]N^{1/2}u.
  \label{eq:mm-commutator-error}
\end{equation}
The reported complete-tree Bessel deficit is
\begin{equation}
  \mathfrak B(h)
  =\left(1-\frac1n\right)\norm{h}^2-\sum_S\norm{K_{M_S}h}^2
  =\sum_S\norm{(I-RR^*)M_SRh}^2.
  \label{eq:mm-bessel-deficit}
\end{equation}
The originating exploration reports rotated examples with negative individual sibling deficits
and Laguerre near-extremizers that consume nearly all descendant slack. Pending persisted
witnesses, these are route-level retractions rather than ledger refutations; they rule out
nodewise positivity and fixed fractional slack allocation within the originating analysis.

\subsection{Schur--Piola transport and local algebra}

For a $1+d$ split write
\[
  H=\begin{pmatrix}h&b^\top\\ b&C\end{pmatrix},\qquad
  v=C^{-1}b,\qquad s=h-b^\top C^{-1}b,\qquad \Delta=\det C,
\]
and $X=\partial_r-v\cdot\nabla_q$. Block algebra, Piola's identity, and Hessian compatibility
formally give
\[
  \operatorname{cof}H
  =\Delta\begin{pmatrix}1&-v^\top\\-v&vv^\top+sC^{-1}\end{pmatrix},
  \qquad
  \operatorname{div}(\Delta X)=0,
\]
\[
  Xv=C^{-1}\nabla_qs,
  \qquad XC=J^\top C=CJ,
  \qquad \Tr(C^{-1}XC)=\operatorname{div}_qv.
\]
If $z=\nabla_q\phi$ is the child target coordinate, then $Xz=0$; for
$x=\nabla\phi$, $Xx_\perp=0$ and $Xx_0=s$. Hence $s^{-1}X$ is the pullback of
$\partial_{x_0}$. In particular, the trace/conformal derivative is constrained by the Schur
connection and is not a free scalar mode.

At a Schur-normal $1+2$ point, let $T\in\operatorname{Sym}_2$, $a,g\in\R^2$, let
$R_\perp e_1=e_2$ and $R_\perp e_2=-e_1$, and use
$\operatorname{sym}(u\otimes v)=\tfrac12(u\otimes v+v\otimes u)$. Put
\[
  B_i=\operatorname{sym}(g\otimes R_\perp^\top e_i).
\]
The fixed-target coordinate form expands exactly as
\begin{align*}
Q(a,T,g)
&=8|a|^2+\frac12\sum_i\norm{\{B_i,T\}}_{\HS}^2-4a\cdot R_\perp Tg\\
&=8\left|a-\frac14R_\perp Tg\right|^2
  +\frac32|Tg|^2+\frac14(\Tr T)^2|g|^2\ge0.
\end{align*}
An independent coordinate expansion checked this identity. What remains open is the invariant
identification of the tensorial lift actually generated by the global operator, and a reduction
covering all higher-dimensional splits. Positivity in $1+2$, $1+3$, and $2+2$ alone would not be
a dimension-free theorem.

\begin{question}[Invariant lift and all-split reduction]
\label{q:mm-invariant-lift}
Derive the multiplier transport in a target-flat Schur--Piola frame without choosing a moving-frame
coefficient by hand. Decompose the resulting Codazzi shape terms into controlled irreducible
pieces for arbitrary split dimensions, and match each negative component with a retained
Letwin/Monge--Amp\`ere square.
\end{question}

\subsection{Two solenoidal channels}

Let $S(z)=\E[C\mid z]$ be the inherited child block of the parent Stein kernel and let $K(z)$ be
the canonical child moment-map Stein kernel. Both solve the same Stein equation, so
\[
  \mathsf D=S-K,
  \qquad \partial_\beta(\rho\mathsf D_{\alpha\beta})=0.
\]
In a two-dimensional child, subject to topology and boundary conditions,
\[
  \rho\mathsf D=R_\perp^\top D^2\lambda R_\perp
\]
is the Airy representation. For a parent datum $h$, the reported canonical flux residual has the
form
\begin{equation}
  r_h=r_{\mathrm{cond}}+(I-\Pi_K)\mathsf D\nabla v.
  \label{eq:mm-hodge-residual}
\end{equation}
Thus a conditional-flux channel and a canonical-versus-inherited Stein-kernel channel must be
controlled separately. The earlier guessed local conserved vector collapsed these two projections
and was withdrawn.

\subsection{The best identified global bottleneck}

The cubic symbol does not cancel in general. With $A_M=H^{-1/2}MH^{-1/2}$, the originating
calculation gives
\[
  C_3(\xi)=2\left[(H^{-1}\xi)^r(\partial_rA_M)(\xi,\xi)
  -(A_M\xi)^r(\partial_rH^{-1})(\xi,\xi)\right].
\]
It is proposed to equal
\[
  2(H^{-1}\xi)^k\xi^\top H^{-1/2}[\Omega_k,M]H^{-1/2}\xi,
  \qquad
  \Omega_k=\frac12\left(H^{-1/2}\partial_kH^{1/2}
  -(\partial_kH^{1/2})H^{-1/2}\right),
\]
exhibiting eigenframe rotation rather than conformal variation. The square-root bridge is the
standard resolvent formula
\[
  [N^{1/2},K_M]
  =\frac1\pi\int_0^\infty t^{1/2}(N+t)^{-1}
  [N,K_M](N+t)^{-1}\dd t,
\]
provided all forms and domains are justified.

\begin{question}[Global square-root commutator]
\label{q:mm-square-root-commutator}
Define a retained nonnegative remainder $\mathfrak R_{\mathrm{Letwin}}(u)$ containing every
available variable-multiplier, Codazzi, Monge--Amp\`ere, corrector, and descendant square, and
prove over the complete Haar tree that
\[
  2\sum_S\operatorname{Re}\langle K_{M_S}Nu,e_S(u)\rangle
  +\sum_S\norm{e_S(u)}^2
  \le \mathfrak B(Nu)+\mathfrak R_{\mathrm{Letwin}}(u)
\]
with a dimension-free constant and stable operator domains. Then show that the complete-tree
estimate supplies the global analytic input required by the fully defined CMH target of
Question~\ref{q:cmh-normalization}.
\end{question}

This is the best identified bottleneck, not yet the sole formal gap. The invariant lift, all-split
reduction, Hodge boundary conditions, one-edge flux identity, and endpoint duality must be closed
alongside it. The originating regression analysis indicates that no scalar, conformal-only, or
nodewise shortcut survives its current models. Persisting those witnesses is still required, but
the resulting proof design is necessarily global and matrix-coupled.
\end{document}
